\section{Related Work}
\label{sec: related_work}

\textbf{Benchmarks for Document Understanding.} A great amount of datasets have emerged to evaluate the DU capabilities of LVLMs. Many datasets focus exclusively on either a single component (\eg table, chart)~\cite{masry-etal-2022-chartqa, zhu2022towards, liu2024visualwebbench, kardas-etal-2020-axcell} or a single page~\cite{Mathew2020DocVQAAD, Mathew2021InfographicVQA} from the full documents. Some recent DU datasets~\cite{tito2023hierarchical, VanLandeghem2023DocumentUD, SlideVQA2023, saadfalcon2023pdftriage, islam2023financebench} attempt to assess multi-page documents, but still exhibit shortcomings in terms of document length (page number), information density (token number) and the construction approaches. Specifically, MP-DocVQA~\cite{tito2023hierarchical} is an extension of DocVQA~\cite{Mathew2020DocVQAAD} and inherently absent of both cross-page and unanswerable questions. Annotating from scratch, DUDE~\cite{VanLandeghem2023DocumentUD} includes a small percentage of cross-page questions (2.1\%) and unanswerable questions (12.7\%). However, due to the relatively short context length (5.3 pages on average) and the use of crowd-sourced annotations, questions in DUDE tend to be less challenging and somewhat less rigorous. SlideVQA features 20-page documents and cross-page questions (12.9\%). Nevertheless, the documents in SlideVQA are in slide-deck format and of relatively low information density. Moreover, these cross-page questions are HotpotQA-style~\cite{yang-etal-2018-hotpotqa} created by instantiating entity graphs and co-referencing in-graph entities across multiple pages. The entity graph from a closed document tends to be sparse and has significant shortcuts (see examples in Appendix~\ref{subsec: Existing Question Editing}). These shortcuts sometimes lead to false cross-page questions that actually do not require answer evidence across different pages. The recent FinanceBench~\cite{islam2023financebench} features both extremely long-context documents and practical, scalable cross-page questions. However, its documents are exclusively financial reports. Additionally, the reference answers are in open-ended formats, making the expert-level manual evaluation indispensable. The above reasons limit the broader applicability of FinanceBench. To our best knowledge, \benchmarkname is the first comprehensive, qualified, and easy-to-use benchmark on the long-context DU task. More detailed descriptions and comparisons are presented in Table~\ref{tab:dataset_comparison}.

\paragraph{Models for Document Understanding.} 
There are two main branches of models for automatic DU tasks. The first approach employs two-stream, OCR-dependent architectures to separately encode textual information (parsed via OCR) and visual information (images and/or layout structures)~\cite{LayoutLMv1, xu-etal-2021-layoutlmv2, LayoutLMv3}. In contrast, the second approach develops OCR-free models that understand documents in an end-to-end manner~\cite{kim2022donut, Pix2Struct}. With the rapid advancement of LVLMs, the latter approach has dominated the current DU solutions. As mentioned above, a range of LVLMs demonstrate promising performance on single-page DU datasets. However, as shown in Section~\ref{sec:evaluation}, even the most advanced LVLMs fall significantly short of achieving satisfactory performance on our benchmark. It reveals that understanding lengthy documents still poses great challenges to current LVLMs.


\paragraph{Long-context LVLMs and LLMs.} 
Lengthy documents necessitate the use of LVLMs or LLMs with extended context sizes. Several benchmarks~\cite{shaham-etal-2022-scrolls, an2023leval, bai2023longbench, zhang2024inftybench} and solutions~\cite{Tworkowski2023FocusedTC, Chen2023LongLoRAEF, Peng2023YaRNEC, bai2024longalign} have been proposed to evaluate and develop long-context LLMs. However, there exists limited related work for long-context LVLMs, leaving this area largely unexplored. Until very recently, contemporary studies~\cite{song2024milebench, jiang2024mantis, lu2024text} assess and/or improve LVLMs' multi-image understanding capabilities. Evaluations on both \benchmarkname and these works indicate that current LVLMs are still not fully equipped to handle long-context DU and many other practical tasks that require extensive contextual comprehension.